\section{Finite-word bounds for a law}\label{sec:lp}

Let $Q=Q_\mu$ for a law on $[0,M]$, and set $Q(0)=0$ for the lower
grid. Choose $k\ge1$ and $r\ge2$. Type $j$ has mark $v_j=Q(j/k)$.
The indices $1,\ldots,k$ remain distinct even when two marks coincide.
For each of the $k^r$ words $w$ of length $r$ in these indices, let
$c_j(w)$ count index $j$, and let $\ell(w)$ be the least span in that
order, evaluated on the marks by~\eqref{eq:recurrence}. Define
\begin{equation}\label{eq:lp}
\begin{split}
 \lambda_{k,r}(\mu)=\min_p\ &\frac1r\sum_w p_w\ell(w),\\
 p_w\ge0,\quad &\sum_w p_w=1,\qquad
 \sum_w p_wc_j(w)=\frac rk\quad(1\le j\le k).
\end{split}
\end{equation}
An individual word can omit entire types. The constraints require balance
only in expectation, and the uniform mixture of the $k$ constant words
is feasible.

\begin{theorem}[Finite-word bounds for a law]\label{thm:profile-words}
For every $k\ge1$ and $r\ge2$,
\[
 \lambda_{k,r}(\mu)-M/k\ \le\ E(\mu)\ \le\ \lambda_{k,r}(\mu)+M/r.
\]
The resulting bracket for $E(\mu)/\pi$ has width $M(1/k+1/r)/\pi$
before numerical evaluation of the word costs.
\end{theorem}
\begin{proof}
Form two lists with $q$ copies of each indexed type, using $Q(j/k)$
for the upper list and $Q((j-1)/k)$ for the lower. Write $e^+$ and $e^-$
for their limiting spans per mark as $q\to\infty$, which exist by
Theorem~\ref{thm:profile-law} applied to their discrete laws. In the
uniform case these type limits were obtained in~\cite{supp} by applying
subadditivity~\cite[Theorem~4.9]{walters} to the span plus one.
The sorted
quantile sample of size $kq$ lies between the two sorted lists. They
share $k-1$ indexed types. Starting with a minimizing lower placement,
delete the $q$ extra zeros and append $q$ copies of the top mark at
spacing $M$, with a gap $M$ to the retained points. This costs at most
$qM$, also when $k=1$ and the retained list is empty. Monotonicity under
decreasing marks now gives
\begin{equation}\label{eq:types}
 e^+-M/k\le e^-\le E(\mu)\le e^+ .
\end{equation}

Take $q$ divisible by $r$ and cut a minimizing upper placement into
$kq/r$ successive words of length $r$. Their empirical distribution
satisfies the balance equations in~\eqref{eq:lp} since each index occurs $q$ times.
Their minimum spans sum to at most the full span. Dividing by $kq$
and sending $q$ to infinity along these multiples gives
$\lambda_{k,r}(\mu)\le e^+$.

For the converse, the feasible polytope in~\eqref{eq:lp} is rational
and bounded. Its real linear objective attains a minimum at a rational
vertex $p$. Clear its denominators with $D$ and take $Dp_w$ copies of
each word. This batch has $Dr/k$ occurrences of each index, an integer
by the balance equations. Lay out repeated batches with every word at
its minimum span and a gap $M$ between successive words. By
Lemma~\ref{lem:closing} every cross-word pair has enough separation.
The mean span of these placements tends to $\lambda_{k,r}(\mu)+M/r$,
so $e^+$ is at most this value. Combine with~\eqref{eq:types}.
\end{proof}

\begin{corollary}[Uniform finite-word bounds]\label{thm:lp}
For the uniform law on $[0,1]$ the types are $j/k$. Writing
$\lambda_{k,r}$ for the value of~\eqref{eq:lp} on these types gives
\begin{equation}\label{eq:effective}
 \frac{\lambda_{k,r}-1/k}{\pi}\le\CC\le
 \frac{\lambda_{k,r}+1/r}{\pi}.
\end{equation}
In particular $\pi\CC=\lim_{k,r\to\infty}\lambda_{k,r}$ jointly.
\end{corollary}

\subsection{Evaluating the word costs}
For example, $k=r=64$ gives width $1/(32\pi)<0.01$ at $M=1$, with $64^{64}$
words in~\eqref{eq:lp}. If directed rational enclosures of all edge
weights $\sqrt{v_i v_j}$ have width at most $\epsilon$, applying
\eqref{eq:recurrence} at both endpoints encloses each word span in
$[\ell^-(w),\ell^+(w)]$ of width at most $(r-1)\epsilon$ since each path has
at most $r-1$ edges. The two rational endpoint programs thus add at most
$(r-1)\epsilon/r$ to the width in mean-span units. They admit exact
vertex enumeration. For a certificate of an upper value, any feasible rational $p$
gives $\sum_w p_w\ell^+(w)/r$. Any rational $y_0,y_j$
satisfying
\[
 y_0+\sum_j y_jc_j(w)\le\ell^-(w)/r\qquad\text{for every word }w
\]
gives the lower value $y_0+(r/k)\sum_j y_j$. Both certificates use
only rational comparisons.

For uniform types, directed square-root and $\pi$ enclosures therefore
compute $\CC$ to arbitrary prescribed accuracy in principle. For a
general law the same procedure requires an effective bound $M$ and
access to arbitrarily accurate enclosures of the indicated quantiles.
Even the sole type at $k=1$ for $\delta_c$ requires access to the atom location
$c$. The theorem gives a structural width for every bounded law.

\subsection{The word LP for square radii}\label{sec:worked-lp-example}

For the profile $f(x)=x^2$, the law $\mu$ of $f(U)$ has quantile
$Q_\mu(t)=t^2$ and mark bound $M=1$, so the bracket applies directly to
the square-radius coefficient $C_2=E(\mu)/\pi$.
This choice also makes the computation exact in integers because the edge between
type indices $i,j$ has length $ij/k^2$.

Take $k=r=6$ and write a word as a string of type indices.
The primal distribution has three words of positive mass, with scaled
spans and probabilities as follows.
\[
 (w,36\ell(w),p_w)\ \in\
 \{(511135,30,1/6),\ (612326,36,1/2),\ (534445,79,1/3)\}.
\]
The probabilities sum to one and the mean count of each index is one.
Taking one, three and two copies, respectively, gives a balanced batch.
A unit gap at each join recovers a feasible line placement.
The spans use all predecessors in \eqref{eq:recurrence}. For instance, in
$511135$ the path through the indices $5,3,5$ already costs $30/36$.
For the dual of \eqref{eq:lp} take
\[
 y_0=0,\qquad 648(y_1,\ldots,y_6)=(-9,-8,13,40,52,60).
\]
For every word, the integer check is
$\sum_j(648y_j)c_j(w)\le 3(36\ell(w))$.
Checking this inequality on all $6^6$ words certifies dual feasibility.
The weighted mean of $\ell(w)/6$ and the dual value $\sum_j y_j$ both
equal $37/162$. Theorem~\ref{thm:profile-words} therefore gives
\[
 \frac{5}{81\pi}\le C_2\le\frac{32}{81\pi}.
\]

\begin{table}[htbp]
\centering
\caption{Square-radius LP bounds. Multiplying the bracket and its width by
$\pi$ gives the rational entries shown. Dashes indicate sizes for which
only the word count and guaranteed width were calculated.}
\label{tab:worked-lp-example}
\begin{tabular}{rrrrcc}
\toprule
$k$ & $r$ & Words $k^r$ & $\lambda_{k,r}(\mu)$ &
Interval for $\pi C_2$ & $\pi$ times width \\
\midrule
4 & 4 & $256$ & $7/32$ & $[-1/32,15/32]$ & $1/2$ \\
6 & 4 & $1296$ & $41/216$ & $[5/216,95/216]$ & $5/12$ \\
6 & 6 & $46656$ & $37/162$ & $[5/81,32/81]$ & $1/3$ \\
16 & 16 & $16^{16}$ & --- & --- & $1/8$ \\
64 & 64 & $64^{64}$ & --- & --- & $1/32$ \\
\bottomrule
\end{tabular}
\end{table}

Table~\ref{tab:worked-lp-example} displays the raw theorem bounds. Its
first lower endpoint is vacuous since spans are nonnegative.
At fixed $r$, finer quantiles reduce the $1/k$ contribution to the scaled
width. At fixed $k$, longer words reduce the $1/r$ contribution.
Even in the tightest computed row, Corollary~\ref{cor:profile-powers}
gives a stronger lower bound and Proposition~\ref{prop:cert-upper} a
stronger upper bound. At $k=r=16$, even one byte per word would require
$2^{64}$ bytes for full enumeration in memory.


\subsection{Power radii}
A second use of the line problem gives an elementary lower bound without
enumerating words.

\begin{corollary}[Adjacent-pair lower bound]\label{cor:profile-adjacent}
For any probability law on $[0,M]$,
\[
 E(\mu)\ge\int_0^1\sqrt{Q_\mu(t)Q_\mu(1-t)}\,dt .
\]
\end{corollary}
\begin{proof}
Let $s_1\le\cdots\le s_n$ be the sorted square roots of the quantile
sample. Each successive gap in a minimizing line order dominates its
adjacent product. Add the product for the last and first marks, at most
$M$. The successor map permutes the square roots, so the rearrangement
inequality bounds the cyclic sum below by $\sum_i s_i s_{n+1-i}$.
Hence $b_n^Q+M$ dominates that sum. After division by $n$, the reverse
pairing tends to the displayed integral by bounded convergence, since
$Q_\mu$ is continuous at both $t$ and $1-t$ for almost every $t$.
\end{proof}

\begin{corollary}[Power radii]\label{cor:profile-powers}
For each fixed real $p>0$ and outer radii $i^p$,
\[
 R_p^*(n)=C_p n^{p+1}+o(n^{p+1}),\qquad
 C_p=E(\operatorname{law}(U^p))/\pi ,
\]
where $U$ is uniform on $(0,1)$. For square radii, $C_2\ge 1/(6\pi)$.
\end{corollary}
\begin{proof}
Take $N_n=n^p$ and $a_{n,i}=(i/n)^p$ in Theorem~\ref{thm:profile-law}.
The quantile is $Q(t)=t^p$. For $p=2$ the integral in
Corollary~\ref{cor:profile-adjacent} is $\int_0^1t(1-t)\,dt=1/6$.
\end{proof}
